\begin{abstract}
%TL;DR (X) A language model that writes build programs makes Minecraft houses from text for cents, and a natural upgrade is to let it imagine the target as an image first.
%TL;DR (Y) Whether the image helps is hard to measure, because a dream loop also adds turns and text, and its extra text is what the builder reads.
%TL;DR (Z) We separate the image from the turns with a compute-matched loop without an image, and the image from its description with caption controls.
%TL;DR (V) The caption of the dream moves multi-view realism; the image alone does not; no arm moves geometry, and all arms build houses players rarely enter and that repeat each other.
A language model that writes short WorldEdit-style build programs turns a text prompt into a
Minecraft house for about four cents, with no training data. A popular upgrade for such agents
is visual imagination: a text-to-image model draws the target, and the agent builds, renders,
compares and repairs until its build matches the picture. Measuring what the picture adds is
harder than it looks, because an imagination loop also adds repair turns and gives the builder
extra text, and either one alone changes the builds. We isolate the image with a no-image loop
that matches the number of turns and the selection rule, and we isolate the information in the
image with a structured caption of it, a caption of a different prompt's image, and a text
specification written without any image. On 133 house prompts, with \gptmini{} as the builder and
\imgmini{} as the image model, the caption arm nearly halves the multi-view render distance to real
houses against the no-image loop (\mvkid{} 0.217 against 0.403, Holm-corrected $p=0.014$), while
the image without a caption ties it. No arm separates from another on geometry, the
no-image loop builds the most enterable houses (68\% against 50\% and 51\%, real 93\%), and in every
arm a fifth to a third of the builds have a near-copy in the same set, against none of the real houses. Visual imagination
buys a program-writing agent appearance, not structure. \TODO{n=256: replace every number in this
abstract with the n=256 card (lands Nov 10) and add the shuffled-caption, text-spec and best-of-4 results.}
% 0.04 USD/build: card run_20261109_121625 cost table (0.040 to 0.048), T51
% MV-KID 0.217 vs 0.403, p_holm 0.009: T51 addendum 2 (n=133); the 2026-11-09 card gives p_holm 0.0136 with more arms
% enterable 0.68 / 0.50 / 0.51, real 0.95: T63
% twin rate 0.21 to 0.32 vs real 0.008 (ratios 26 to 40): card run_20261109_121625 novelty table, real from T91 (n=256)
\end{abstract}
